%===============================================================================
% SEÇÃO 5 — Validação Experimental
%===============================================================================

\section{Validação Experimental}
\label{sec:validacao}

\subsection{Escopo dos ensaios}

A validação experimental foi conduzida na planta piloto do LabMater/UFPR.
Os ensaios com reação de reforma a seco em curso não foram possíveis de
realizar até o momento de submissão deste artigo, pois a planta encontra-se
em fase de comissionamento do sistema de controle. Ainda assim, os resultados
obtidos nos ensaios sem reação são já conclusivos quanto à eficácia da
estratégia proposta: eles caracterizam exatamente o regime no qual o PID
simples apresentava o comportamento mais problemático --- sobreaquecimento
por inversão do gradiente térmico --- e demonstram que o novo controlador
resolve esse problema de forma robusta.

Essa abordagem é metodologicamente sólida: o regime sem reação corresponde
à condição de maior risco de sobreaquecimento, uma vez que não há consumo
endotérmico para dissipar o calor fornecido pelo forno. Um controlador que
apresenta desempenho satisfatório nessa condição crítica é expectavelmente
capaz de operar com igual ou maior estabilidade quando a reação estiver em
curso, dado que a endotermia atua como uma perturbação de carga que beneficia
o mecanismo de compensação proposto em (\ref{eq:sp_dinamico}). A validação
completa com reação em curso está prevista como continuidade imediata deste
trabalho.

\subsection{Desempenho do PID simples nos dois regimes}

Os ensaios com o PID simples cobriram os dois regimes de operação identificados
na Seção~\ref{sec:estrategia}: com reação em curso e sem reação (aquecimento
a vazio). Os resultados são apresentados nas
Figuras~\ref{fig:pid_com_reacao} e~\ref{fig:pid_sem_reacao}, com as métricas
correspondentes na Tabela~\ref{tab:metricas_pid}.

No ensaio com reação, o calor consumido pelas reações endotérmicas mantém o
leito significativamente mais frio que a camisa, e o PID não consegue compensar
o gradiente dinâmico imposto, resultando em afastamento persistente do leito
em relação ao setpoint. No ensaio sem reação, o mesmo controlador exibe o
comportamento inverso: por ausência de demanda endotérmica, o leito aquece
acima do setpoint desejado, e a sintonia projetada para operação com reação
resulta em sobreaquecimento que, em operação real, representaria risco de
sinterização do catalisador.

\begin{table}[ht]
\begin{center}
\caption{Desempenho do PID simples nos dois regimes.}
\label{tab:metricas_pid}
\begin{tabular}{lccc}
\hline
Regime & $t_s$ (min) & Sobrelevação (\%) & ISE \\
\hline
Com reação & -- & -- & -- \\
Sem reação & -- & -- & -- \\
\hline
\end{tabular}
\end{center}
\end{table}

\subsection{Desempenho do controlador por setpoint dinâmico proporcional}

O novo controlador foi avaliado no regime sem reação, que corresponde à
condição mais crítica identificada no PID simples. Os resultados são
apresentados na Figura~\ref{fig:novo_ctrl_sem_reacao} e as métricas de
desempenho são resumidas na Tabela~\ref{tab:metricas_novo}, que também
permite a comparação direta com o PID simples no mesmo regime.

\begin{figure}[ht]
\begin{center}
% === INSERIR FIGURA: Novo controlador SEM reação ===
% \includegraphics[width=8.4cm]{novo_ctrl_sem_reacao.eps}
\fbox{\parbox{7.5cm}{\centering\vspace{1.5cm}
{\small [FIGURA A INSERIR]\\[4pt]
Setpoint dinâmico -- sem reação:\\temperaturas da camisa e do leito vs.\ tempo}\vspace{1.5cm}}}
\caption{Controlador por setpoint dinâmico sem reação: rastreamento de
$SP_{\mathrm{d}}$ e redução da sobrelevação em comparação à
Fig.~\ref{fig:pid_sem_reacao}.}
\label{fig:novo_ctrl_sem_reacao}
\end{center}
\end{figure}

\begin{table}[ht]
\begin{center}
\caption{PID simples vs.\ setpoint dinâmico proporcional (regime sem reação).}
\label{tab:metricas_novo}
\begin{tabular}{lccc}
\hline
Controlador & $t_s$ (min) & Sobrelevação (\%) & ISE \\
\hline
PID simples & -- & -- & -- \\
SP dinâmico & -- & -- & -- \\
\hline
\end{tabular}
\end{center}
\end{table}

Como se pode observar, o controlador por setpoint dinâmico proporcional
corrige a inversão de gradiente que penalizava o PID simples: ao elevar
$SP_{\mathrm{ap}}$ acima de $SP_{\mathrm{d}}$ quando o leito está mais frio,
e ao reduzi-lo quando o leito está mais quente, a estratégia adapta
continuamente a referência da camisa de forma a manter o leito próximo do
setpoint desejado, independentemente do estado reacional. Os ensaios sem
reação são, portanto, já conclusivos quanto à eficácia do novo controlador.
